% year: תשפ"ג
% type: מבחן
% semester: א
% moed: א
% number: 2ב
% points: 10
% categories: ivt
% source: exams/מבחנים/תשפג/מועד א תשפג.pdf

\begin{question}
תהי $f(x)$ פונקציה רציפה על כל הישר, ונניח כי לכל $x$ מתקיים $|f(x)|\le 7$. הוכיחו כי למשוואה $2x+f(x)=3$ יש לפחות פתרון אחד.
\end{question}

\begin{hint}
הגדירו $g(x)=2x+f(x)-3$ ומצאו נקודה שבה $g$ שלילית ונקודה שבה $g$ חיובית, בעזרת החסם $|f(x)|\le 7$.
\end{hint}

\begin{solution}
נגדיר $g(x)=2x+f(x)-3$. הפונקציה $g$ רציפה על כל $\R$ כסכום של פונקציות רציפות. מהנתון $-7\le f(x)\le 7$ לכל $x$, ולכן:
\[ g(-10)=-20+f(-10)-3\le -23+7=-16<0,\qquad g(10)=20+f(10)-3\ge 17-7=10>0. \]
$g$ רציפה בקטע הסגור $[-10,10]$ ומקבלת בקצותיו ערכים בעלי סימנים מנוגדים, ולכן לפי \textbf{משפט ערך הביניים} (משפט בולצאנו) קיימת $c\in(-10,10)$ שבה $g(c)=0$, כלומר $2c+f(c)=3$. לכן למשוואה יש לפחות פתרון אחד. $\blacksquare$
\end{solution}
